% Generated draft: requires mathematical and source review.
\documentclass[11pt]{article}
\usepackage[margin=25mm]{geometry}
\usepackage{fontspec}
\setmainfont{Latin Modern Roman}
\usepackage{amsmath,amssymb,mathtools,unicode-math}
\setmathfont{Latin Modern Math}
\usepackage{enumitem,needspace,xurl,hyperref}
\setlength{\parindent}{0pt}
\setlength{\parskip}{5pt}
\setlength{\emergencystretch}{3em}
\newcommand{\N}{\mathbb{N}}
\newcommand{\Z}{\mathbb{Z}}
\newcommand{\Q}{\mathbb{Q}}
\newcommand{\R}{\mathbb{R}}
\newcommand{\C}{\mathbb{C}}
\newcommand{\F}{\mathbb{F}}
\newcommand{\A}{\mathbb{A}}
\newcommand{\PP}{\mathbb{P}}
\newcommand{\E}{\mathbb{E}}
\newcommand{\eps}{\varepsilon}
\newcommand{\dd}{\,\mathrm{d}}
\begin{document}
\section*{UnsolvedMath ID: 11000139}
\textbf{Source problem code:} catalogue entry
\subsection{Definitions and mathematical statement}
Let $F$ be a surface in the source's hyperbolic ideal-cell-decomposition setting, with Teichm\"uller space $\mathcal T(F)$. For a marked hyperbolic structure $X$, the source associates an ideal cell decomposition $\Delta(X)$ and projective simplicial coordinates $[s_e(X)]$ on its edges. The dual fatgraph is the graph with one vertex per cell and one edge per shared side, with cyclic orders induced by the oriented surface. Use the projective simplicial coordinates as Strebel length coordinates on this fatgraph to construct a marked conformal structure $\Phi(X)$. This describes the intended map
\[
 X\longmapsto(\Delta(X),[s_e(X)])\longmapsto\Phi(X)\in\mathcal T(F).
\]
The excerpt does not specify the peripheral decorations needed for the hyperbolic construction, normalization of projective coordinates, or zero-edge/cell-boundary conventions. Those choices must be supplied before the map is completely defined.

\textbf{Problem 9 (Bounded Distortion Conjecture).} The resulting self-map of $\mathcal T(F)$ has bounded distortion in the Teichm\"uller metric. A natural comparison interpretation is uniformly bounded displacement:
\[
 \exists C(F)<\infty\quad\forall X\in\mathcal T(F),\qquad
 d_{\mathrm{Teich}}(X,\Phi(X))\le C(F).
\]
With $d_{\mathrm{Teich}}(X,Y)=\tfrac12\inf\log K(f)$ over marking-compatible quasiconformal maps, this would uniformly control the quasiconformal comparison of the hyperbolic-derived conformal structure and the Strebel-derived one. The source does not define ``bounded distortion'' explicitly; a Lipschitz/quasi-isometry bound comparing pairs of points would be a different assertion. The displayed displacement formulation is therefore an interpretation requiring source confirmation, not a certified replacement for the phrase.

The appended historical problem, attributed to Ed Witten in the early 1990s, asks for an orbifold compactification of $\mathcal M(F)$ with cells described by generalized fatgraphs, permitting matrix-model calculations such as [16,19,22]. The source allows the desired compactification to be Deligne--Mumford or another one. It says its preceding combinatorial compactification is instead a more general stratified space and announces a following solution guided by Dennis Sullivan. That solution is absent from this record; the compactification announcement must not be confused with a solution of the Bounded Distortion Conjecture.
\subsection{Short English statement}
Compare the conformal structure from a hyperbolic surface with the one obtained by using its simplicial coordinates as Strebel fatgraph lengths, and ask for uniform Teichm\"uller distortion control. The source separately announces a solution to an orbifold-compactification problem.
\subsection{Sources}
\begin{itemize}
\item Local UnsolvedMath registry, permanent ID 11000139, source code catalogue entry; no separate proposer or proposal year supplied.
\item Farb, edited volume \emph{Problems on Mapping Class Groups and Related Topics}, Problem 9, PDF page 120; \url{https://www.math.uchicago.edu/~farb/papers/mcgbook.pdf}. The registry records PDF-font-structure extraction and access on 2026-07-29; its inventory distinguishes 329 labeled headings by source page.
\item The text attributes the compactification challenge to Ed Witten in the early 1990s and guidance for its announced solution to Dennis Sullivan. The registry summary associates the source author with Penner; no additional precise citation for the announced solution is supplied.
\item Source labels [16,19,22] refer to calculations motivating matrix models, without full bibliography entries in the record.
\end{itemize}
\end{document}
